\documentclass[conference]{IEEEtran}
\IEEEoverridecommandlockouts
\usepackage{cite}
\usepackage{amsmath,amssymb,amsfonts}
\usepackage{algorithmic}
\usepackage{graphicx}
\usepackage{textcomp}
\usepackage{xcolor}
\usepackage{url}
% Figures: upload fig1_pipeline_architecture.png through fig6_nlp_metrics.png to Overleaf root
\def\BibTeX{{\rm B\kern-.05em{\sc i\kern-.025em b}\kern-.08em
    T\kern-.1667em\lower.7ex\hbox{E}\kern-.125emX}}
\begin{document}

\title{Transfer Learning for Dental Pathology Detection in Panoramic Radiographs: YOLOv11 vs. Zero-Shot Foundation Models}

\author{\IEEEauthorblockN{1\textsuperscript{st} Ernane Ferreira Rocha Junior}
\IEEEauthorblockA{\textit{Dept. of Computer Engineering and Automation} \\
\textit{Federal University of Rio Grande do Norte (UFRN)}\\
Natal, RN, Brazil \\
ernane.ferreira.100@ufrn.edu.br}
\and
\IEEEauthorblockN{2\textsuperscript{nd} Ignacio S\'{a}nchez-Gendriz}
\IEEEauthorblockA{\textit{Northeast Strategic Technologies Center} \\
\textit{CETENE / MCTI}\\
Recife, PE, Brazil\\
ignacio.gendriz@cetene.gov.br}
\and
\IEEEauthorblockN{3\textsuperscript{rd} Luiz Affonso Guedes}
\IEEEauthorblockA{\textit{Dept. of Computer Engineering and Automation} \\
\textit{Federal University of Rio Grande do Norte (UFRN)}\\
Natal, RN, Brazil\\
affonso@dca.ufrn.br}
}

\maketitle

\begin{abstract}
Panoramic dental radiography is the most common supplementary imaging exam in dental practice, yet automated detection systems for it require annotated training data that is scarce and expensive to obtain. Zero-shot vision-language models offer an annotation-free path, but how accurately they localize pathologies on domain-specific images is not well characterized. This work compares Grounding DINO, a state-of-the-art zero-shot object detector, against a YOLOv11m model fine-tuned on 678 images from the DentexChallenge 2023 public dataset, targeting three pathologies: Caries, Periapical lesion, and Impacted tooth. On the DentexChallenge test split ($n = 103$), Grounding DINO achieved mAP@50 = 0.00 for all three classes, while the fine-tuned YOLOv11m reached mAP@50 = 0.499 (mAP@50:95 = 0.321), with Impacted tooth at 0.566 and Caries at 0.303. A three-stage pipeline was then applied to 50 private OPGs: YOLOv11 detection (466 total detections, mean 9.3 per image), spontaneous recall evaluation against dentist-written clinical descriptions (mean SR = 88.3\% across 30 evaluable images), and LLM-based structured pre-clinical report generation via Gemini 2.5 Flash (BERTScore F1 = 0.780 on 50 reports). The results show that zero-shot detection is insufficient for clinical OPG analysis at standard IoU thresholds, and that supervised fine-tuning on a small public dental dataset produces a usable detection baseline.
\end{abstract}

\begin{IEEEkeywords}
panoramic radiography, dental pathology detection, YOLOv11, Grounding DINO, zero-shot detection, transfer learning, DentexChallenge, automated report generation
\end{IEEEkeywords}

\section{Introduction}

Panoramic dental radiography (orthopantomography, OPG) produces a single wide-field exposure covering the entire dentition, mandible, maxilla, and surrounding structures. It is the standard first-line imaging request in dental clinical practice, used for screening caries, periapical lesions, bone loss, impacted teeth, retained roots, and orthodontic findings \cite{ludlow2006}. Systematic interpretation takes time and requires experience, and access to dental radiologists varies considerably across healthcare settings.

Automated detection systems based on convolutional neural networks and, more recently, transformer architectures have shown promising results on specific dental imaging tasks. Tuzoff et al. \cite{tuzoff2019} applied Faster R-CNN to tooth detection and numbering in OPGs. Ekert et al. \cite{ekert2019} used a U-Net for periapical lesion detection. Alqhtani et al. \cite{alqhtani2025} benchmarked multiple YOLO variants against 17 fine-grained dental anomalies, reporting mAP@50 values between 0.54 and 0.78 depending on pathology. All these approaches depend on large, fully annotated bounding box datasets, which take considerable clinical expertise to produce.

A possible path around that bottleneck comes from vision-language models (VLMs) trained at internet scale. These models learn open-vocabulary visual representations that can be queried at inference time using text prompts, without any task-specific fine-tuning \cite{radford2021}. Grounding DINO \cite{liu2024gdino}, presented at ECCV 2024, extends the DINO transformer architecture with grounded pre-training and achieves state-of-the-art zero-shot detection across standard benchmarks (52.5 AP on COCO zero-shot). Deploying such a model directly on OPGs, prompted with disease names and without annotation effort, is an appealing idea.

This work tests whether that idea holds in practice. Specifically: at a standard IoU threshold of 0.5, does Grounding DINO detect Caries, Periapical lesion, and Impacted tooth at a useful level of precision? If not, what does supervised fine-tuning on the largest available public dental dataset achieve on the same task, and with what training resources?

Grounding DINO achieved mAP@50 = 0.00 on all three classes. YOLOv11m fine-tuned on 678 DentexChallenge images reached mAP@50 = 0.499 overall, with per-class results that reflect training data imbalance directly. A complete end-to-end pipeline then applied those detections to 50 private clinical OPGs and fed the output to a multimodal LLM to generate structured pre-clinical reports, evaluated against dentist-written reference descriptions using standard NLP metrics.

\section{Related Work}

\subsection{Supervised Dental Pathology Detection}

Convolutional detectors applied to OPGs generally report mAP in the 0.50--0.80 range for well-represented classes. Alqhtani et al. \cite{alqhtani2025} found that performance varies sharply by pathology: visually distinctive findings such as impactions and large restorations outperform subtle ones such as early caries and periodontal bone loss by 20--30 mAP points within the same model. This per-class variability appears in our results as well. The DentexChallenge 2023 \cite{dentex2023}, a MICCAI workshop challenge, provided a public benchmark with disease-level bounding box annotations across three pathology classes, and we use it as both training source and evaluation benchmark.

\subsection{Zero-Shot and Open-Vocabulary Detection}

Grounding DINO \cite{liu2024gdino} uses a Swin Transformer image encoder paired with a BERT-based text encoder, fused through a feature enhancer, to produce open-vocabulary bounding box predictions from text prompts. Prior work applying CLIP \cite{radford2021} to zero-shot classification in medical imaging found that general-purpose vision-language representations transfer to clinical domains for coarse classification, but noted performance gaps relative to supervised models \cite{zhang2024clip}. Localization at IoU $\geq$ 0.5 is substantially harder than classification; the gap between zero-shot and supervised performance is larger. Our experiments confirm this.

\subsection{Automated Radiology Report Generation}

Large multimodal language models have been applied to chest radiograph reporting \cite{liu2025report}, with evaluation using BLEU \cite{papineni2002}, ROUGE-L \cite{lin2004}, and BERTScore \cite{zhang2020bert}. BERTScore F1, computed from contextual token embeddings, correlates more reliably with human quality ratings than n-gram metrics when reference and hypothesis texts differ substantially in length and style \cite{zhang2020bert}. That is the case in our Stage~3 evaluation. No comparable study exists for panoramic dental radiograph reporting.

\section{Materials and Methods}

\subsection{Datasets}

\textit{DentexChallenge 2023 (training and evaluation).}
The DentexChallenge 2023 quadrant-enumeration-disease subset \cite{dentex2023} is the largest publicly available OPG dataset with instance-level disease annotations (CC0 license). It contains approximately 678 images with at least one bounding box annotation from four disease categories: Caries, Deep Caries (merged with Caries in preprocessing), Periapical Lesion, and Impacted tooth. The class distribution is skewed: Caries annotations number in the thousands, while Periapical lesion has 158 annotated instances. Images were split 70/15/15 into train ($n = 474$), validation ($n = 101$), and test ($n = 103$) sets using stratified random sampling (seed = 42). Augmented copies of rare-class training images were added to address class imbalance, resulting in 950 effective training images.

\textit{Private OPG dataset (pipeline evaluation).}
Fifty panoramic radiograph images were collected from clinical practice (JPEG format, mean resolution approximately $2976 \times 1488$ pixels). Each image has a corresponding written clinical description produced by a trained dentist. Thirteen examinations additionally have an audio recording of the verbal report. Bounding box annotations are not available for this dataset; it was used exclusively for the Stage~2 recall and Stage~3 report generation evaluations.

\subsection{Stage 1: Detection Model}

\textit{Preprocessing and augmentation.}
Images were resized to $640 \times 640$ pixels using letterbox padding. Training augmentation included CLAHE (clip limit 3.0) for X-ray contrast normalization, horizontal flip ($p = 0.5$), random brightness and contrast perturbation ($\pm 0.2$, $p = 0.5$), Gaussian noise ($p = 0.3$), and affine transformations (scale 0.9--1.1, translate $\pm 5\%$, rotate $\pm 5°$, $p = 0.4$), implemented with Albumentations. Color augmentations (HSV hue and saturation) were disabled because OPGs are grayscale. Images containing Periapical lesion or Impacted tooth annotations received two additional augmented copies each, tripling their effective training frequency.

\textit{YOLOv11m fine-tuning.}
YOLOv11m (medium variant, approximately 20M parameters) was initialized from COCO-pretrained weights and fine-tuned on the DentexChallenge training split using the Ultralytics framework \cite{jocher2024}. Hyperparameters: AdamW optimizer, initial learning rate $10^{-3}$, linear decay to 1\% over training, batch size 8, image size $640 \times 640$, up to 100 epochs with early stopping (patience = 20). Training ran on an Apple M5 MacBook Air (16~GB unified memory) using the Metal Performance Shaders backend. The model reached best validation mAP@50 = 0.547 at epoch 27 and stopped at epoch 41.

\textit{Grounding DINO zero-shot baseline.}
Grounding DINO Tiny (\texttt{IDEA-Research/grounding-dino-tiny}, Hugging Face Transformers) was evaluated on the same test split without fine-tuning. The three class names were used as text prompts. Detection thresholds: box confidence $\geq 0.25$, text score $\geq 0.20$. Predictions were matched to ground-truth boxes at IoU $\geq 0.5$, consistent with COCO evaluation.

\begin{figure*}[t]
\centerline{\includegraphics[width=\textwidth]{fig1_pipeline_architecture}}
\caption{Three-stage pipeline: YOLOv11 detection (Stage~1), spontaneous recall evaluation against dentist descriptions (Stage~2), and LLM-based report generation via Gemini 2.5 Flash (Stage~3).}
\label{fig:pipeline}
\end{figure*}

\subsection{Stage 2: Spontaneous Recall Evaluation}

The 50 private OPGs were processed with the fine-tuned YOLOv11m. For each image, the corresponding dentist description was parsed to identify mentions of the three YOLO-trained classes through substring matching. The Spontaneous Recall (SR) score for image $i$ is:

\begin{equation}
SR(i) = \frac{|\hat{C}(i) \cap C^{*}(i)|}{|C^{*}(i)|}
\label{eq:sr}
\end{equation}

where $\hat{C}(i)$ is the set of classes detected by YOLOv11 and $C^{*}(i)$ is the set of classes mentioned in the dentist description, both restricted to the three trained classes. The detector processes each image with a fixed confidence threshold of 0.25 and receives no information from the reference text. Images for which no YOLO-trained class appeared in the description were excluded from SR statistics (20 images), leaving 30 evaluable images.

\subsection{Stage 3: Automated Report Generation}

Structured pre-clinical reports were generated for all 50 private OPGs using Gemini 2.5 Flash via the OpenRouter API. Each request contained: (1) the OPG image in base64 encoding, and (2) a summary of detections from Stage~1 (class, confidence, and spatial location for the top 15 detections). The reference clinical description was not included in the prompt. The model was instructed to produce a five-section structured report using FDI tooth notation and to report only findings visible in the image.

Generated reports were compared to dentist-written reference descriptions using BLEU-4 \cite{papineni2002}, ROUGE-L \cite{lin2004}, and BERTScore F1 \cite{zhang2020bert} with RoBERTa-large embeddings.

\section{Results}

\subsection{Zero-Shot vs. Fine-Tuned Detection}

Table~\ref{tab:detection} shows detection results on the DentexChallenge 2023 test split ($n = 103$ images), the only evaluation set with ground-truth bounding boxes.

\begin{table}[htbp]
\caption{Detection performance on DentexChallenge 2023 test set (IoU = 0.5, conf $\geq$ 0.25)}
\begin{center}
\small
\begin{tabular}{lcccc}
\hline
\textbf{Model} & \textbf{Car.} & \textbf{Peri.} & \textbf{Imp.} & \textbf{mAP@50} \\
\hline
GDINO (zero-shot)     & 0.000 & 0.000 & 0.000 & 0.000 \\
YOLOv11m (fine-tuned) & 0.303 & 0.095 & 0.566 & 0.499 \\
\hline
\multicolumn{5}{l}{\scriptsize Car.~=~Caries; Peri.~=~Periapical; Imp.~=~Impacted.} \\
\multicolumn{5}{l}{\scriptsize YOLOv11m: mAP@50:95 = 0.321, P = 0.583, R = 0.550.}
\end{tabular}
\label{tab:detection}
\end{center}
\end{table}

Grounding DINO produced bounding box outputs when queried with the three class names, but none reached IoU $\geq$ 0.5 with any ground-truth box across the test split. The model generates text-anchored candidate regions that do not align closely enough with annotated pathology boundaries for clinical localization.

The fine-tuned YOLOv11m reached mAP@50 = 0.499, with Precision = 0.583 and Recall = 0.550. Per-class results track the training data distribution: Impacted tooth (AP@50 = 0.566) has a distinctive radiographic appearance and reasonable training coverage; Caries (AP@50 = 0.303) is the most prevalent pathology but also the most morphologically variable; Periapical lesion (AP@50 = 0.095) had only 158 annotated training instances.

\begin{figure}[htbp]
\centerline{\includegraphics[width=\columnwidth]{fig2_training_curves}}
\caption{Validation mAP@50 over training epochs for YOLOv11n and YOLOv11m on the DentexChallenge validation split. YOLOv11m peaked at 0.547 (epoch 27) before early stopping at epoch 41.}
\label{fig:training}
\end{figure}

\begin{figure}[htbp]
\centerline{\includegraphics[width=\columnwidth]{fig3_model_comparison}}
\caption{Per-class AP@50 for Grounding DINO (zero-shot) and YOLOv11m (fine-tuned) on the DentexChallenge test set.}
\label{fig:comparison}
\end{figure}

\subsection{Detection on Private OPGs (Stage 1)}

Applied to the 50 private OPGs, the fine-tuned YOLOv11m produced 466 total detections (mean $9.3 \pm 3.2$ per image, range 3--20). By class: Caries 387 (83.0\%), Impacted tooth 59 (12.7\%), Periapical lesion 20 (4.3\%). The Caries proportion is consistent with its prevalence in routine clinical OPGs; the low Periapical count follows from the model's weaker AP for that class.

Fig.~\ref{fig:det_example} shows a representative Stage~1 output on a private OPG. Multiple Caries detections appear bilaterally in the posterior sectors, and a single Impacted tooth is flagged in the upper right third molar region. No Periapical lesion was detected in this image, consistent with that class's lower recall. The color scheme follows the application convention: red for disease-category findings (Caries, Periapical lesion) and green for tooth-status findings (Impacted tooth).

\begin{figure*}[t]
% Upload: results/figures/stage1/det_11.jpg → Overleaf root as det_11.jpg
\centerline{\includegraphics[width=\textwidth]{det_11}}
\caption{Representative Stage~1 detection output on a private OPG. Red bounding boxes indicate Caries detections (bilateral posterior sectors); the green bounding box marks an Impacted tooth in the upper right third molar region. Confidence scores are displayed on each label. No Periapical lesion was detected in this image.}
\label{fig:det_example}
\end{figure*}

\begin{figure}[htbp]
\centerline{\includegraphics[width=\columnwidth]{fig4_detection_frequency}}
\caption{Class distribution of the 466 detections produced on the 50 private OPGs.}
\label{fig:frequency}
\end{figure}

\subsection{Spontaneous Recall (Stage 2)}

Of the 50 private OPGs, 30 had dentist descriptions mentioning at least one YOLO-trained class. Mean SR across those 30 images was 88.3\% (SD = 30.8\%; median = 1.000). Table~\ref{tab:sr} shows the distribution.

\begin{table}[htbp]
\caption{Stage 2 spontaneous recall distribution ($n = 30$ evaluable images)}
\begin{center}
\begin{tabular}{lcc}
\hline
\textbf{SR band} & \textbf{$n$} & \textbf{\%} \\
\hline
SR = 0.0                           & 3  & 10.0 \\
$0 < \text{SR} < 0.5$              & 0  &  0.0 \\
$0.5 \leq \text{SR} < 0.8$        & 1  &  3.3 \\
SR $\geq$ 0.8                      & 26 & 86.7 \\
\hline
\end{tabular}
\label{tab:sr}
\end{center}
\end{table}

The three images with SR = 0.0 all had Periapical lesion as the only YOLO-testable class mentioned in the description, consistent with the model's low AP@50 = 0.095 for that class.

Fig.~\ref{fig:val_example} illustrates a Stage~2 validation overlay for an image with SR = 1.0. The dentist's description mentioned an impacted tooth, and the model detected exactly that class in the lower right third molar area. Green boxes denote classes whose detection matches the reference description; orange boxes would indicate detections outside the expected set for that image.

\begin{figure}[htbp]
% Upload: results/figures/stage2/val_20.jpg → Overleaf root as val_20.jpg
\centerline{\includegraphics[width=\columnwidth]{val_20}}
\caption{Stage~2 spontaneous recall overlay (SR = 1.0). The green bounding box marks an Impacted tooth detection in the lower right third molar region, matching the class mentioned in the corresponding dentist description. The green tint covers the full image area where the detection was confirmed.}
\label{fig:val_example}
\end{figure}

\begin{figure}[htbp]
\centerline{\includegraphics[width=\columnwidth]{fig5_spontaneous_recall}}
\caption{Spontaneous recall distribution across the 30 evaluable private OPGs.}
\label{fig:recall}
\end{figure}

\subsection{Report Generation Quality (Stage 3)}

Gemini 2.5 Flash generated structured five-section pre-clinical reports for all 50 images without error (136,257 input tokens, 26,057 output tokens). Table~\ref{tab:nlp} shows NLP metrics against dentist-written reference descriptions.

\begin{table}[htbp]
\caption{NLP evaluation of generated reports vs. dentist reference descriptions ($n = 50$)}
\begin{center}
\begin{tabular}{lccc}
\hline
\textbf{Metric} & \textbf{Mean} & \textbf{SD} & \textbf{Embeddings} \\
\hline
BLEU-4        & 0.0011 & 0.0012 & n-gram              \\
ROUGE-L       & 0.0319 & 0.0174 & n-gram              \\
BERTScore F1  & 0.7803 & 0.0135 & RoBERTa-large       \\
\hline
\end{tabular}
\label{tab:nlp}
\end{center}
\end{table}

Low BLEU-4 and ROUGE-L values are expected given that generated reports average approximately 300 words across five structured sections while dentist reference descriptions average approximately 45 words of free text. Exact n-gram overlap between these two formats is low regardless of semantic accuracy. BERTScore F1 = 0.780 indicates reasonable semantic alignment at the level of contextual word embeddings.

A concrete example illustrates this gap: for Image\_05, the generated report identified impacted teeth 18 and 38 with correct angulation descriptions, listed multiple carious lesions in FDI notation with severity estimates, and recommended appropriate follow-up actions. The dentist's reference note for the same image read: ``Grossly carious 48 and deep caries in 36 with pulp exposure. Septic root of 25.'' BLEU and ROUGE scores for this pair are near zero because the texts differ in structure and length; BERTScore F1 for this pair was 0.791, reflecting that both texts describe the same pathological findings.

\begin{figure}[htbp]
\centerline{\includegraphics[width=\columnwidth]{fig6_nlp_metrics}}
\caption{NLP metric distributions across 50 generated reports. BLEU-4 and ROUGE-L are near zero due to structural differences between generated and reference texts; BERTScore F1 captures semantic agreement.}
\label{fig:nlp}
\end{figure}

\section{Discussion}

\subsection{Localization Failure in Zero-Shot Detection}

Grounding DINO is not blind to dental anatomy: when prompted with class names, it generates candidate boxes in plausible image regions. The failure is at localization precision, not semantic recognition. Standard mAP evaluation counts a detection as correct only if $\text{IoU}(\hat{b},\, b^{*}) \geq 0.5$. For small lesions like periapical radiolucencies (typically a few millimeters in diameter at clinical scale), a box that roughly surrounds the affected region but is 30\% too large fails this threshold entirely. Grounding DINO was trained to ground language to visual regions from natural image datasets, where objects are much larger and more salient than millimeter-scale dental pathologies. That spatial precision is simply not in what the model learned.

In practice, this means zero-shot VLMs are not usable as detection components in any pipeline that requires IoU-based evaluation at clinical thresholds. They may serve well for rough localization, region proposals, or qualitative visualization, but not for quantitative clinical assessment as formulated here.

\subsection{Transfer Learning from a Small Public Dataset}

The mAP@50 = 0.499 from 678 training images (950 after augmentation) is meaningful for two concrete reasons. First, it confirms that the transfer from ImageNet and COCO pre-training generalizes to the X-ray domain at this dataset size when the augmentation matches the modality: CLAHE, disabled color perturbation, and oversampling of rare classes. Second, Impacted tooth at AP@50 = 0.566 is a clinically actionable result. Third molar impactions are among the most commonly missed findings in routine OPG review, and a detection system at this performance level can function as a second reader.

\subsection{Periapical Lesion: A Data Problem}

Periapical lesion AP@50 = 0.095 follows from 158 annotated training examples. Periapical radiolucencies are small (2--10~mm), radiographically subtle in early stages, and easily confused with anatomical structures or noise artifacts. Improving this class would require more training data, targeted augmentation, or semi-supervised learning from unlabeled OPGs. This is a data limitation, not an architecture one, and reporting it directly is more useful than obscuring it.

\subsection{Web Prototype}

To assess whether the pipeline is compatible with clinical workflows in terms of latency and interface, a web application was built on top of the three-stage pipeline. The backend runs on FastAPI and exposes a REST API that receives an uploaded OPG, runs Stage~1 detection via the YOLOv11m model, feeds the result to the Stage~3 Gemini 2.5 Flash call, and stores the annotated image and generated report. The frontend is a React single-page application with four main components: an annotation canvas that overlays detection bounding boxes on the radiograph with a toggle to compare the raw and annotated views; a structured findings list grouped by class; a report viewer with tabbed sections mirroring the five-section report structure; and a voice input component that allows the clinician to record corrections or additional observations directly into the interface.

End-to-end latency from image upload to a fully rendered annotated radiograph and structured pre-clinical report was approximately 15 seconds on an Apple M5 MacBook Air with a standard broadband connection to the OpenRouter API. The dominant cost is the Gemini API call, not the local YOLO inference, which completes in under one second.

Fig.~\ref{fig:web_output} shows the annotated output produced by the system for one of the private OPGs processed through the web interface. Detection boxes are rendered on the stored annotated image and returned as a JPEG alongside the structured report.

\begin{figure}[htbp]
% Upload: web/backend/storage/annotated/ba491db79c114d5093c4d1e8300b7804_annotated.jpg
%         → Overleaf root as web_annotated.jpg
% Alternatively replace with a full UI screenshot (web_screenshot.png)
\centerline{\includegraphics[width=\columnwidth]{web_annotated}}
\caption{Annotated OPG output generated by the web prototype. Red bounding boxes indicate Caries detections across the posterior dentition, as returned by the Stage~1 YOLOv11m model and stored by the FastAPI backend. A structured pre-clinical report is displayed alongside this image in the frontend interface.}
\label{fig:web_output}
\end{figure}

The prototype confirms that the computational requirements are compatible with commodity hardware: no GPU is required at inference time on the server side, and the Apple MPS backend handles YOLO inference locally if the pipeline runs offline. A production deployment would require server-side inference and appropriate patient data handling under applicable regulations, which are outside the scope of this work.

\subsection{Limitations}

Four constraints limit the scope of these conclusions. First, the evaluation uses a single public benchmark (DentexChallenge 2023); results on datasets from different equipment, populations, or annotation protocols may differ. Second, Stage~2 SR and Stage~3 NLP metrics are not substitutes for clinical accuracy evaluation because the private dataset has no bounding box ground truth. Third, the generated reports were not reviewed by clinicians; BERTScore correlation with dental-report quality has not been established for this domain. Fourth, training and inference ran on an Apple M5 MPS backend; throughput figures may differ on other hardware, though the architecture itself is GPU-agnostic.

\section{Conclusion}

Grounding DINO, tested without fine-tuning, achieved mAP@50 = 0.00 on the DentexChallenge 2023 test set for all three target classes because its localization precision is insufficient at IoU $\geq$ 0.5 for millimeter-scale dental pathologies. YOLOv11m fine-tuned on 678 annotated images from the same public dataset reached mAP@50 = 0.499 overall, with Impacted tooth at 0.566, Caries at 0.303, and Periapical lesion at 0.095. Applied to 50 private clinical OPGs, the pipeline achieved a mean spontaneous recall of 88.3\% against dentist descriptions and produced structured pre-clinical reports with BERTScore F1 = 0.780.

Foundation models with zero-shot object detection are not yet suitable for clinical dental detection at standard evaluation thresholds. Supervised fine-tuning on a small public dataset, with augmentation tailored to the modality, produces a viable detection baseline. A React/FastAPI web prototype of the complete pipeline delivers a structured pre-clinical report from image upload in approximately 15 seconds on commodity hardware, with no GPU required at inference time.

Future work will focus on scaling the annotated training set, exploring semi-supervised approaches for Periapical lesion, extending detection to the full 31-class clinical vocabulary, and obtaining clinical validation of the generated reports by specialist dentists.

\section*{Acknowledgment}

The authors thank the DentexChallenge 2023 organizers for making the dataset publicly available under a CC0 license. The private OPG dataset was collected under institutional protocols for anonymized clinical data use.

\bibliographystyle{IEEEtran}
\bibliography{references}

\end{document}
